\documentclass[11pt,a4paper]{article}

\usepackage[T1]{fontenc}
\usepackage[utf8]{inputenc}
\usepackage{lmodern}
\usepackage[ngerman]{babel}
\usepackage[a4paper,margin=2.2cm]{geometry}
\usepackage{array}
\usepackage{longtable}
\usepackage{tabularx}
\usepackage{xcolor}
\usepackage{hyperref}

\setlength{\parindent}{0pt}
\setlength{\parskip}{0.5em}
\renewcommand{\arraystretch}{1.35}

\newcommand{\answerline}[1]{\rule{#1}{0.4pt}}
\newcommand{\term}[1]{\fcolorbox{black!20}{blue!6}{\strut #1}}

\title{\textbf{Aufgabenblatt: Eigenschaften verteilter Systeme}}
\author{Modul 321 -- Lektion 2}
\date{}

\begin{document}

\maketitle

\textbf{Name:} \answerline{7cm} \hfill \textbf{Datum:} \answerline{3.5cm}

\vspace{0.8em}

\fcolorbox{blue!40}{blue!6}{
  \parbox{\dimexpr\linewidth-2\fboxsep-2\fboxrule-6pt\relax}{
    \textbf{Ziel der Lektion:} Sie können typische Eigenschaften verteilter Systeme
    an realistischen Szenarien erklären und geeignete Architekturentscheidungen
    für Retry, Idempotenz, Resilience sowie Transaktionsgrenzen begründen.
  }
}

\section*{Hinweise}
\begin{itemize}
  \item Antworten Sie in ganzen Sätzen oder präzisen Stichworten.
  \item Begründen Sie fachliche Entscheidungen knapp, aber klar.
  \item Verwenden Sie bei Bedarf kleine Skizzen oder Ablaufdiagramme.
\end{itemize}

\section*{1. Begriffe sauber erklären}
Erklären Sie die folgenden Begriffe in \textbf{1 bis 3 eigenen Sätzen}.

\begin{longtable}{|>{\raggedright\arraybackslash}p{0.28\linewidth}|>{\raggedright\arraybackslash}p{0.64\linewidth}|}
\hline
\textbf{Begriff} & \textbf{Ihre Erklärung} \\
\hline
\term{Idempotenz} & \\[2.4cm] \hline
\term{Retry} & \\[2.0cm] \hline
\term{Resilience / Design for Failure} & \\[2.4cm] \hline
\term{Event Sourcing} & \\[2.2cm] \hline
\end{longtable}

\section*{2. Retry beurteilen}
Kreuzen Sie an oder notieren Sie \textbf{Ja / Nein / Nur unter Bedingungen}. Begründen Sie anschliessend jede Entscheidung in \textbf{einem kurzen Satz}.

\begin{tabularx}{\linewidth}{|>{\raggedright\arraybackslash}p{0.52\linewidth}|>{\centering\arraybackslash}p{0.14\linewidth}|>{\raggedright\arraybackslash}X|}
\hline
\textbf{Situation} & \textbf{Retry?} & \textbf{Begründung} \\
\hline
Ein Request endet wegen Timeout, die Ursache ist vermutlich ein temporärer Netzwerkfehler. & & \\[1.0cm] \hline
Ein Request schlägt wegen ungültiger Kundendaten fehl. & & \\[1.0cm] \hline
Eine Zahlung wurde möglicherweise schon verarbeitet, aber die Antwort kam nicht beim Client an. & & \\[1.2cm] \hline
Ein externer Dienst liefert HTTP~503 wegen Überlastung. & & \\[1.0cm] \hline
\end{tabularx}

\section*{3. Idempotenz im Praxisfall}
Ein Shop sendet an den Payment-Service:

\begin{quote}
\texttt{POST /payments} mit Betrag, Bestellnummer und Kreditkartendaten.
\end{quote}

Nach 4 Sekunden erhält der Shop keine Antwort mehr und sendet dieselbe Anfrage erneut.

\begin{enumerate}
  \item Welches fachliche Risiko entsteht?
  \item Warum reicht ein einfacher Retry hier nicht aus?
  \item Beschreiben Sie, wie ein \term{Idempotency Key} das Problem entschärfen kann.
  \item Welche Information müsste der Service zu diesem Schlüssel mindestens speichern?
\end{enumerate}

\section*{4. Resilience auf Architektur-Ebene}
Ein Online-Shop besteht aus folgenden Teilen:

\begin{itemize}
  \item Frontend
  \item Produktkatalog
  \item Warenkorb
  \item Empfehlungssystem
  \item Payment-Service
\end{itemize}

Das Empfehlungssystem fällt aus.

\begin{enumerate}
  \item Beschreiben Sie ein \textbf{schlechtes Verhalten} des Gesamtsystems.
  \item Beschreiben Sie ein \textbf{resilientes Verhalten} des Gesamtsystems.
  \item Nennen Sie \textbf{zwei technische Massnahmen}, die ein solches resilientes Verhalten unterstützen.
  \item Warum ist \emph{degradierter Betrieb} oft besser als ein Totalausfall?
\end{enumerate}

\section*{5. Transaktionsgrenzen verstehen}
Betrachten Sie folgendes Szenario:

\begin{quote}
Ein Bestellprozess umfasst drei Services: \textbf{Order-Service}, \textbf{Payment-Service} und \textbf{Inventory-Service}. Der Order-Service erstellt die Bestellung, der Payment-Service autorisiert die Zahlung und der Inventory-Service reserviert den Bestand.
\end{quote}

Bearbeiten Sie die Aufgaben:
\begin{enumerate}
  \item Warum ist eine einzige klassische Datenbanktransaktion über alle drei Services schwierig?
  \item Welche Probleme entstehen, wenn ein Service langsam ist oder ausfällt?
  \item Was bedeutet in diesem Kontext \term{eventual consistency} in eigenen Worten?
  \item Nennen Sie ein Beispiel, wie man fachlich mit einem Teilfehler umgehen könnte.
\end{enumerate}

\section*{7. Event Sourcing einordnen}
Ordnen Sie die folgenden Aussagen \textbf{zu} und kommentieren Sie kurz, ob sie zu \term{Event Sourcing} passen.

\begin{enumerate}
  \item Es wird nur der aktuelle Zustand gespeichert, frühere Änderungen interessieren nicht.
  \item Ereignisse wie \texttt{OrderPlaced} oder \texttt{PaymentAuthorized} werden dauerhaft gespeichert.
  \item Ein Read-Model kann aus einer Folge von Events neu aufgebaut werden.
  \item Event Sourcing ist immer die beste Wahl für einfache CRUD-Anwendungen.
\end{enumerate}

\section*{8. Entwurfsaufgabe}
Sie entwerfen eine kleine Plattform zur \textbf{Online-Reservation von Arbeitsplätzen} in einem Schulhaus.

Die Plattform soll:
\begin{itemize}
  \item Reservierungen annehmen,
  \item Verfügbarkeiten anzeigen,
  \item Bestätigungs-E-Mails senden,
  \item mit temporären Störungen einzelner Komponenten umgehen können.
\end{itemize}

Bearbeiten Sie dazu:
\begin{enumerate}
  \item Zeichnen oder beschreiben Sie eine einfache Architektur mit mindestens vier Komponenten.
  \item Nennen Sie eine Stelle, an der \term{Retry} sinnvoll sein kann.
  \item Nennen Sie eine Stelle, an der \term{Idempotenz} wichtig sein kann.
  \item Beschreiben Sie ein Beispiel für \term{degradierten Betrieb}.
  \item Begründen Sie, ob Sie in diesem Szenario \term{Event Sourcing} einsetzen würden oder bewusst nicht.
\end{enumerate}

\end{document}
